\section{Natural Images as Boundary Evidence}
\label{sec:cifar}
The local second-moment representation in Section~\ref{sec:task-geometry} is an architecture identity for these experiments. Its population risk and movement bounds do not, however, predict natural-image behavior until their task moments, normal update budgets, and stability constants have been measured. The existing CIFAR-10 and DTD results delimit this gap; they are not a validation of the new finite-horizon bounds. Full protocols, capacity scans, DC interventions, and local derivative diagnostics are retained in Appendix~\ref{app:natural-details}.

\begin{table}[t]
\centering\small
\caption{Natural-image boundary evidence: validation OLS accuracy (\%), initialization $\to$ final budget. Three development seeds per row; CIFAR frozen readouts share the same encoders. Rank-8 fitting uses training labels. DTD uses official partition 1.}
\label{tab:natural-boundary}
\setlength{\tabcolsep}{4pt}
\begin{tabular}{llrr}
\toprule
Dataset&Representation&Initial&Final\\\midrule
CIFAR&Encoder&36.31&36.04\\
&Original eight-dimensional head&23.49&23.71\\
&Supervised rank-8 readout&35.75&35.59\\\midrule
DTD&Encoder&8.65&8.51\\
&Zero-DC encoder&8.78&8.56\\
\bottomrule
\end{tabular}
\end{table}

CIFAR uses grayscale images, 64 local $11\times11$ filters, squaring, GAP, and an eight-dimensional head. At a fixed frozen encoder, supervised rank-8 least squares recovers almost all of the full encoder's validation MSE advantage over the original head. The approximately twelve-point accuracy gap is already present at initialization (Table~\ref{tab:natural-boundary}). It demonstrates available class-dependent compression, not a large training-induced deletion of class information. Reduced-rank fitting optimizes training MSE, not validation accuracy. The separate five-seed final-test comparison has encoder accuracy near $36.5\%$ across rate ratios; equal-energy validation comparisons do not establish a stable speed-independent gain.

On DTD, objective improvement concentrates encoder variance in an existing image statistic associated with DC response. Removing that response at evaluation or training with spatially zero-mean filters does not yield clear encoder class gains. Zero-DC training instead emphasizes a statistic associated with pixel variance. These outcomes are consistent with objective--task mismatch, but neither spectral concentration nor removal of one dominant statistic identifies the cause of weak class learning.

The existing label-conditioned CIFAR derivative diagnostic has eligible opposing contributions in only four of 27 state--stream combinations. All four reverse the first-step risk sign, with changes of order $10^{-9}$; none preserves both predicted signs at twenty updates. Predictions and the first update share the same gradient and use validation labels. This is a local calculation check, not an independent finite-horizon test. Source-state task coupling, separate future-gradient errors, and usable curvature bounds have not been measured for these runs. We therefore do not claim that the new theory explains their small magnitude or short persistence. Testing that claim requires calibration states or batches separated from the future update stream, rather than further searches for classification gains.
